\section{重回帰分析}
\label{s10_MultipleRegression}

\subsection{授業内容}
\subsubsection{科目の中でのこのコマの位置づけ}
\begin{tabular}[t]{ccccccc}
	記述統計[3] & 確率[1] & モデル[4] & 推測・推定[0] & 検定[0] & 実践[2]
\end{tabular}

\subsubsection{概要}
回帰分析によって，データにモデルを当てはめて係数の推定値を得るということを行った。データにもっとも適した形を求めて計算を行ったが，そのこととモデルが適していたかどうかは別である。この点を理解するためには，モデルの適合度について考える必要がある。モデルが適合しているということを踏まえた上で，説明する変数が増える重回帰分析へと議論は展開する。ここでは説明変数の読み取り方への注意を促すべく，偏相関係数の考え方を解説し，偏回帰係数，標準化偏回帰係数などの用語について理解する。最後に，多重共線性の問題についても指摘する。
\subsubsection{コマ主題細目(3-5項目箇条書き)}
\begin{description}


	\item[重回帰分析]説明変数が複数に増えた場合の回帰分析は，重回帰分析Multiple Regression Analysisと呼ぶ。数式的な表現としては項が増えるだけであるが，データのプロットを考えると三次元以上の回帰平面に拡張されていること，それでも面は湾曲していない，すなわち線形関係であることを理解する。
	      \begin{flushright}
		      $\to$ \textcite{kosugi2019}のP.168,図10.7の回帰平面図を参照
	      \end{flushright}

	\item[部分相関と偏相関] 続く偏回帰係数の理解に先立って，部分相関の考え方を理解しておく必要がある。回帰分析によって，被説明変数の分散が説明される部分と残差の部分に分離させられること，残差と説明変数は相関しないことを確認し，残差は説明変数の影響を除外したものと考えることができる。これは条件を必ずしも統制できない事象において，統計モデル的に統制していることでもある。これを踏まえて，変数間関係を改めて理解する。


	\item[偏回帰係数]部分相関の説明において，残差が分離されることを見てきたが，残差同士の相関関係はとくに偏相関ということができる。これは共通する変数からの影響を除いた変数間関係であり，これを見るだけでも変数間関係を見ることにつながる。とくに残差変数を使っての回帰分析を行うとき，得られる係数が偏回帰係数である。重回帰分析における回帰係数はこのように条件付きで考えなければならない。また変数を増やすことで重相関係数の増加することを理解する。
	      \begin{flushright}
		      $\to$\textcite{kosugi2018}のP.60--63.
	      \end{flushright}

	\item[標準化された係数]回帰係数については，素点による回帰係数が理解しやすいが，単位に左右されるものでもあるので相対比較をする場合はすべての変数を標準化した標準化解を用いた方がわかりやすい。標準化解を用いても適合度は変化しないことなど留意点とともに有用性を理解する。
	      \begin{flushright}
		      $\to$\textcite{kosugi2018}のP.66--67.
	      \end{flushright}


	      % \item[多重共線性] 複数の変数による予測を行う場合，説明変数同士が相関していると推定値が不正確になることがある。これをとくに多重共線性の問題と呼ぶ。
	      % \begin{flushright}
	      %     $\to$多重共線性の問題については，\textcite{Shimizu201705}のP.86--89。数学的には\textcite{Haebara200206}のP.244--247を参考にすると良い。計算機科学的には\textcite{Matsuura201610}のP.103--104の見方もおもしろい。
	      % \end{flushright}


\end{description}
\subsubsection{キーワード}
\begin{itemize}
	\item 残差
	\item 決定係数
	\item 重回帰分析
	\item 偏相関係数，偏回帰係数
	\item 標準化解
\end{itemize}
\subsection{授業情報}
\paragraph{コマの展開方法}
講義
\paragraph{標準シラバスにおける位置づけ}
科目番号5 心理学統計法；1.心理学で用いられる統計手法；C 記述統計：回帰
\subsubsection{予習・復習課題}
\paragraph{予習}
説明変数が1つの回帰分析について，数式的表現や最適化関数(最小二乗法)の意味について概念的な理解で十分なので再確認しておくこと。後者について今回も同じ基準が用いられるが，説明変数が複数になる場合は解析的算出には行列計算の知識が必要である。行列計算については，二年次の統計に関する講義で扱う予定であるが，興味があるものは「線形代数」をキーワードに予習すると良い。初学者には\textcite{Yuki201810}が，あるいは\textcite{Okada200804}がよい。
\paragraph{復習}
\textcite{Toyoda201706}のP.89--98は本時の内容がわかりやすくまとめられているので，通読しておくと良い復習になる。また重回帰分析は広く用いられる手法なので，\textcite{kosugi2016a}の第二章などを読んで，どのような使い方をされているのか，また前提となっている条件や利用上の問題点などについて，講義時間内で触れられなかった諸問題についての知識を補完しておいてもらいたい。
